\documentclass{article}
\usepackage{amsmath}
\usepackage{amssymb}
\usepackage{amsthm}
\usepackage{mathtools}
\usepackage{geometry}
\usepackage{hyperref}

\newcommand{\rk}{\operatorname{rank}}
\newcommand{\GL}{\operatorname{GL}}
\newcommand{\R}{\mathbb{R}}
\newcommand{\Mat}{\mathrm{Mat}}
\newcommand{\Stab}{\operatorname{Stab}}

\theoremstyle{definition}
\newtheorem{theorem}{Theorem}[section]
\newtheorem{proposition}[theorem]{Proposition}
\newtheorem{lemma}[theorem]{Lemma}
\newtheorem{corollary}[theorem]{Corollary}
\newtheorem{definition}[theorem]{Definition}
\newtheorem{remark}[theorem]{Remark}

\title{Task 0: Validate the Statement of Conjecture B}
\date{}

\begin{document}
\maketitle

\section{Goal}

The purpose of this task is to determine whether the current quotient statement in Conjecture B is mathematically well-defined.

The current candidate statement is:
\[
\mathcal C_{d,S}^r / G_d^{W_0}
\cong
\Gamma_{d,S}^r / G_d^{W_0},
\qquad
G_d^{W_0}:=\{g\in G_d : g_0W_0g_L^{-1}=W_0\}.
\]

\section{Main check}

We must determine whether the critical locus
\[
\mathcal C_{d,S}^r
=
\{W\in \Gamma_{d,S}^r : \mu(W)=P_SW^\ast\}
\]
is stable under the action of $G_d^{W_0}$.

\begin{remark}
Since
\[
\mu(g\cdot W)=g_L\mu(W)g_0^{-1},
\]
the critical locus can only be $G_d^{W_0}$-stable if the distinguished product $P_SW^\ast$ is itself fixed by all elements of $G_d^{W_0}$.
\end{remark}

\section{Decision}

\begin{lemma}[How $G_d^{W_0}$ moves the critical locus]
Let $g\in G_d^{W_0}$. Then
\[
g\cdot \mathcal C_{d,S}^r
=
\left\{
W'\in \Gamma_{d,S}^r
\;\middle|\;
\mu(W')=g_L(P_SW^\ast)g_0^{-1}
\right\}.
\]
\end{lemma}

\begin{proof}
By the equivariance of the product map,
\[
\mu(g\cdot W)=g_L\mu(W)g_0^{-1}.
\]
By the stability of $\Gamma_{d,S}^r$ under $G_d^{W_0}$, proved in the group-action note, one also has
\[
W\in \Gamma_{d,S}^r \Longrightarrow g\cdot W\in \Gamma_{d,S}^r.
\]
Therefore, if $W\in \mathcal C_{d,S}^r$, then $g\cdot W\in \Gamma_{d,S}^r$ and
\[
\mu(g\cdot W)=g_L(P_SW^\ast)g_0^{-1}.
\]
This proves the claimed description of the translated locus.
\end{proof}

\begin{proposition}[Non-stability for positive rank]\label{prop:nonstable-positive-rank}
Assume $r>0$. Then $\mathcal C_{d,S}^r$ is \emph{not} stable under the action of $G_d^{W_0}$.
\end{proposition}

\begin{proof}
Choose a scalar $\lambda\in \mathbb R^\times$ with $\lambda\neq 1$, and define
\[
g_0=I_{d_0},
\qquad
g_\ell=I_{d_\ell}\ \text{ for }1\le \ell\le L-1,
\qquad
g_L=\lambda P_S+P_Q.
\]
Since $P_S$ and $P_Q$ are complementary projectors, $g_L$ is invertible with
\[
g_L^{-1}=\lambda^{-1}P_S+P_Q.
\]
Moreover,
\[
P_Qg_L^{-1}=P_Q,
\]
so
\[
g_0W_0g_L^{-1}
=
\Sigma_{XY}P_Q(\lambda^{-1}P_S+P_Q)
=
\Sigma_{XY}P_Q
=
W_0.
\]
Thus $g\in G_d^{W_0}$.

Now let $W\in \mathcal C_{d,S}^r$. Then
\[
\mu(g\cdot W)=g_L(P_SW^\ast)g_0^{-1}
=
(\lambda P_S+P_Q)P_SW^\ast
=
\lambda P_SW^\ast.
\]
Because $r>0$, the matrix $P_SW^\ast$ has rank $r$ and is therefore nonzero. Since $\lambda\neq 1$, we obtain
\[
\mu(g\cdot W)\neq P_SW^\ast.
\]
Hence $g\cdot W\notin \mathcal C_{d,S}^r$, proving that $\mathcal C_{d,S}^r$ is not $G_d^{W_0}$-stable.
\end{proof}

\begin{remark}[Edge case $r=0$]
If $r=0$, then $P_S=0$ and $P_SW^\ast=0$. In that special case the condition
\[
\mu(W)=0
\]
is preserved by every endpoint action, so $\mathcal C_{d,S}^0$ is stable under $G_d^{W_0}$. Thus the failure of stability is a genuinely positive-rank phenomenon.
\end{remark}

\begin{definition}[Correct stabilizer subgroup]
Define
\[
H_S
:=
\Stab_{G_d^{W_0}}(P_SW^\ast)
=
\left\{
g\in G_d^{W_0}
\;\middle|\;
g_L(P_SW^\ast)g_0^{-1}=P_SW^\ast
\right\}.
\]
\end{definition}

\begin{corollary}[Task 0 conclusion]
For $r>0$, the quotient
\[
\mathcal C_{d,S}^r/G_d^{W_0}
\]
is not well-defined as an orbit space, because $\mathcal C_{d,S}^r$ is not stable under the action of $G_d^{W_0}$.

The correct replacement is to formulate Conjecture B in two parts:
\begin{enumerate}
\item every $G_d^{W_0}$-orbit in $\Gamma_{d,S}^r$ meets $\mathcal C_{d,S}^r$;
\item for every $W\in \mathcal C_{d,S}^r$,
\[
(G_d^{W_0}\cdot W)\cap \mathcal C_{d,S}^r
=
H_S\cdot W.
\]
\end{enumerate}
Equivalently, once the existence statement is proved, one obtains a bijection
\[
\Gamma_{d,S}^r/G_d^{W_0}
\cong
\mathcal C_{d,S}^r/H_S.
\]
\end{corollary}

\section{Expected outcomes}

\subsection*{Case 1: the critical locus is stable}

If
\[
g_L(P_SW^\ast)g_0^{-1}=P_SW^\ast
\qquad
\text{for all } g\in G_d^{W_0},
\]
then the current quotient statement survives unchanged.

\subsection*{Case 2: the critical locus is not stable}

If not, then the corrected statement should be
\[
\Gamma_{d,S}^r/G_d^{W_0}
\cong
\mathcal C_{d,S}^r/H_S,
\qquad
H_S:=\Stab_{G_d^{W_0}}(P_SW^\ast),
\]
or equivalently:
each $G_d^{W_0}$-orbit in $\Gamma_{d,S}^r$ meets $\mathcal C_{d,S}^r$ in exactly one $H_S$-orbit.

\begin{remark}
Task 0 shows that, under the standing positive-rank setup of the conjecture, we are in Case 2.
Case 1 survives only in the degenerate rank-zero situation.
\end{remark}

\section{Deliverables}

This task is complete once the file contains:

\begin{itemize}
\item a lemma deciding stability or non-stability of $\mathcal C_{d,S}^r$ under $G_d^{W_0}$;
\item the final corrected statement of Conjecture B;
\item the definition of the stabilizer subgroup $H_S$ if needed.
\end{itemize}

\end{document}
